Distributed service-oriented applications~\cite{channabasavaiah2003migrating} power everything from global cloud infrastructure to real-time industrial automation.  They operate over networks of computers where services are the primary abstraction. For example, in edge AI systems, WiFi and cellular base stations may provide storage and AI inference/training services to end devices, while also offloading compute-intensive tasks to cloud services. As another example, a UAV network may adopt a service-oriented design to coordinate a ground station and multiple drones: each drone exposes flight-control services invoked by the ground station, while drones may also request services from one another (e.g., for sensing or data processing).  Service frameworks provide the software infrastructure to develop, deploy, and manage such applications.  

NDNSF is motivated by applications in which provider availability and task suitability change over time. A UAV's position and sensing conditions may change, while an inference provider's workload and available memory may vary. Earlier service information describes an offered capability, but may not indicate whether the provider can serve a particular request at that moment. Applications therefore need request-specific participation feedback before selecting executors. NDNSF connects this feedback to user selection and subsequent execution; feedback is not a resource reservation and may become outdated before execution.

In today’s Internet, gRPC~\cite{grpc} is widely used for high-performance remote procedure calls, REST~\cite{fielding2000architectural} supports stateless HTTP-based interactions, and message-queue systems~\cite{john2017mqsurvey} enable asynchronous decoupled communication. Despite their widespread adoption, particular deployments of these frameworks can face limitations in dynamic and large-scale environments. \emph{Scalability} can be constrained when provider selection requires repeated point-to-point exchanges. \emph{Reliability} can be impacted by endpoint binding and load concentration, which create performance bottlenecks and potential single points of failure. \emph{Security} based only on channels protects data in transit; object-level confidentiality requires additional encryption.
% since IP does not natively bind security to named content or enforce fine-grained access control over data objects.

Named Data Networking (NDN)~\cite{zhang2014named} provides an alternative, data-centric networking model in which communication is driven by named data rather than host addresses. Consumers issue Interests to retrieve Data packets, and each Data packet is cryptographically signed to ensure integrity and authenticity, with optional encryption for confidentiality. NDN improves \emph{scalability} through in-network caching and native multicast delivery, allowing multiple consumers to retrieve the same published Data. It enhances \emph{reliability} by decoupling communication from fixed endpoints, enabling retrieval from caches or alternative sources under failures or mobility. Signatures validated under appropriate trust rules support authenticity and integrity; encryption is additionally needed to restrict content access. These properties make NDN a strong foundation for secure service-oriented applications.

\emph{Despite the above advantages, building service-oriented applications over NDN remains challenging because applications must coordinate discovery, access control, and invocation.} To address this gap, we design the Named Data Networking Service Framework (NDNSF), a modular framework for building secure and distributed service-oriented applications with minimal NDN-specific complexity. NDNSF provides three core capabilities: (i) \textbf{service publication and secure discovery}, (ii) \textbf{fine-grained, name-based access control using attribute-based encryption}, and (iii) \textbf{service invocation with dynamic provider selection}. In particular, NDNSF introduces a \textbf{one-to-many service model with customizable provider-selection strategies}, enabling applications to choose providers based on runtime metrics such as computational capacity, queue state, and network delay.

Providers assess task requirements against their current resources and return positive ACKs when willing to participate. The user then applies a customizable selection strategy to choose among these candidates.

We evaluate NDNSF against NDN Service Call (NSC)~\cite{meirovitch2021nsc} and gRPC. In our single-provider wired experiments with packet loss, NDNSF achieves higher completion rates than NSC and rates comparable to gRPC. In scenarios with multiple heterogeneous and intermittently unavailable providers in both wired and wireless networks, NDNSF achieves higher success rates in the evaluated configurations by selecting among available providers, whereas the configured baselines use a fixed provider. We also evaluate selective acknowledgment under provider overload and user-defined provider-selection strategies.

NDNSF's basic service invocation workflow is driven by four Data types and combines protected request dissemination, provider participation feedback, user selection, and execution. Its authorization design has three aspects:
\begin{enumerate}
\item \textbf{Confidential One-to-Many Service Discovery.} NAC-ABE encrypts discovery descriptors containing task requirements that providers need to decide whether to participate, without requiring the requester to list provider identities. 
\item \textbf{Aggregated Service Credentials.} Under one ABE authority and the same set of public parameters, each entity can hold one decryption key (DKEY) encoding multiple service permissions. Section~\ref{sec:credential-cost} separates credential counts from distribution and update costs.
\item \textbf{Reusable Identity Credentials with Simple Trust Rules.} ABE policies represent service permissions, so baseline trust rules can focus on which identities and keys may sign messages rather than listing each entity's service permissions. Identity-signing credentials remain reusable across applications sharing the same trust configuration.
\end{enumerate}
These features use the access-control capabilities of NAC-ABE~\cite{zhang2018nac,dulal2023nac-abe}. Permission aggregation is inherited from ABE; we analyze its organization for service permissions.
